Um die Kontinuitätsgleichung für die Wahrscheinlichkeitsdichte zu zeigen, beginnen wir mit der Definition der Wahrscheinlichkeitsdichte 

\[
\rho(\mathbf{r}, t) = \psi^*(\mathbf{r}, t) \psi(\mathbf{r}, t).
\]

Die zeitliche Ableitung der Wahrscheinlichkeitsdichte ist gegeben durch

\[
\frac{\partial}{\partial t} \rho(\mathbf{r}, t) = \frac{\partial}{\partial t} (\psi^* \psi) = \psi^* \frac{\partial \psi}{\partial t} + \psi \frac{\partial \psi^*}{\partial t}.
\]

Verwenden wir die Schrödinger-Gleichung und ihre komplex konjugierte Form:

\[
\mathrm{i} \hbar \frac{\partial \psi}{\partial t} = \left(-\frac{\hbar^2}{2m} \Delta + V(\mathbf{r})\right) \psi,
\]

\[
-\mathrm{i} \hbar \frac{\partial \psi^*}{\partial t} = \left(-\frac{\hbar^2}{2m} \Delta + V(\mathbf{r})\right) \psi^*,
\]

erhalten wir für die zeitliche Ableitung der Wahrscheinlichkeitsdichte

\[
\frac{\partial}{\partial t} \rho(\mathbf{r}, t) = \frac{1}{\mathrm{i} \hbar} \left( \psi^* \left(-\frac{\hbar^2}{2m} \Delta \psi + V \psi \right) - \psi \left(-\frac{\hbar^2}{2m} \Delta \psi^* + V \psi^* \right) \right).
\]

Die Terme mit dem Potential \( V \) heben sich gegenseitig auf, sodass wir

\[
\frac{\partial}{\partial t} \rho(\mathbf{r}, t) = -\frac{\hbar}{2 \mathrm{i} m} \left( \psi^* \Delta \psi - \psi \Delta \psi^* \right).
\]

Nun betrachten wir den Ausdruck für den Wahrscheinlichkeitsstrom

\[
\mathbf{j}(\mathbf{r}, t) = \frac{\hbar}{2 \mathrm{i} m} \left( \psi^* \nabla \psi - \psi \nabla \psi^* \right).
\]

Die Divergenz des Wahrscheinlichkeitsstroms ist

\[
\nabla \cdot \mathbf{j}(\mathbf{r}, t) = \frac{\hbar}{2 \mathrm{i} m} \left( \nabla \cdot (\psi^* \nabla \psi) - \nabla \cdot (\psi \nabla \psi^*) \right).
\]

Verwenden wir die Produktregel, erhalten wir

\[
\nabla \cdot (\psi^* \nabla \psi) = \nabla \psi^* \cdot \nabla \psi + \psi^* \Delta \psi,
\]

\[
\nabla \cdot (\psi \nabla \psi^*) = \nabla \psi \cdot \nabla \psi^* + \psi \Delta \psi^*.
\]

Daraus folgt

\[
\nabla \cdot \mathbf{j}(\mathbf{r}, t) = \frac{\hbar}{2 \mathrm{i} m} \left( \psi^* \Delta \psi - \psi \Delta \psi^* \right).
\]

Setzen wir dies in die Kontinuitätsgleichung ein, erhalten wir

\[
\frac{\partial}{\partial t} \rho(\mathbf{r}, t) + \nabla \cdot \mathbf{j}(\mathbf{r}, t) = -\frac{\hbar}{2 \mathrm{i} m} \left( \psi^* \Delta \psi - \psi \Delta \psi^* \right) + \frac{\hbar}{2 \mathrm{i} m} \left( \psi^* \Delta \psi - \psi \Delta \psi^* \right) = 0.
\]

Damit ist die Kontinuitätsgleichung gezeigt.